% Generated draft: requires mathematical and source review.
\documentclass[11pt]{article}
\usepackage[margin=25mm]{geometry}
\usepackage{fontspec}
\setmainfont{Latin Modern Roman}
\usepackage{amsmath,amssymb,mathtools,unicode-math}
\setmathfont{Latin Modern Math}
\usepackage{enumitem,needspace,xurl,hyperref}
\setlength{\parindent}{0pt}
\setlength{\parskip}{5pt}
\setlength{\emergencystretch}{3em}
\newcommand{\N}{\mathbb{N}}
\newcommand{\Z}{\mathbb{Z}}
\newcommand{\Q}{\mathbb{Q}}
\newcommand{\R}{\mathbb{R}}
\newcommand{\C}{\mathbb{C}}
\newcommand{\F}{\mathbb{F}}
\newcommand{\A}{\mathbb{A}}
\newcommand{\PP}{\mathbb{P}}
\newcommand{\E}{\mathbb{E}}
\newcommand{\eps}{\varepsilon}
\newcommand{\dd}{\,\mathrm{d}}
\begin{document}
\section*{UnsolvedMath ID: 30003985}
\textbf{Source problem code:} catalogue entry
\subsection{Definitions and mathematical statement}
The supplied record states a five-part research programme for random walks. It does not specify the acting group or semigroup, the class of spaces, or the hypotheses under which the requested conclusions are sought. Consequently, it does not determine a single universally quantified conjecture.

The following notation explains a conventional interpretation without supplying the missing hypotheses. Let \(S\) denote the unspecified acting group or semigroup, let \(X\) denote the unspecified space, and let \(\mu\) be a probability measure on \(S\). Whenever the topology, measurable structures, and action make these expressions meaningful, the usual interpretation of the semigroup closure is
\[
\Gamma=\overline{\langle\operatorname{supp}\mu\rangle},
\]
where \(\langle\operatorname{supp}\mu\rangle\) is the semigroup generated by the support of \(\mu\), and the closure is taken in the topology of \(S\). For \(x_0\in X\), its orbit is \(\Gamma x_0=\{g x_0:g\in\Gamma\}\), whereas its orbit closure is \(\overline{\Gamma x_0}\), with closure in \(X\). These sets need not coincide; the source calls \(\Gamma x_0\) an orbit closure without resolving this distinction.

Under a left-action convention, independent increments \(G_1,G_2,\ldots\) with common law \(\mu\) give positions
\[
Z_0=x_0,\qquad Z_k=G_k\cdots G_1x_0\quad(k\geq1).
\]
One possible meaning of the empirical measures in the record is
\[
\nu_n=\frac1n\sum_{k=0}^{n-1}\delta_{Z_k}\quad(n\geq1),
\]
where \(\delta_x\) is the probability measure concentrated at \(x\). The record does not define \(\nu_n\); in particular, it does not distinguish these random measures from their expectations \(\frac1n\sum_{k=0}^{n-1}\operatorname{Law}(Z_k)\).

A Borel probability measure \(\eta\) on \(X\) is \(\mu\)-stationary when
\[
\eta(A)=\int_S\eta(g^{-1}A)\,\mathrm d\mu(g)
\]
for every Borel set \(A\subseteq X\), provided the integral is defined; here \(g^{-1}A=\{x\in X:gx\in A\}\) denotes a preimage. Write \(\mathcal P_\mu(X)\) for the set of such measures. A standard interpretation of an ergodic stationary measure is an extreme point of this convex set: if
\[
\eta=t\eta_1+(1-t)\eta_2,\qquad 0<t<1,\qquad \eta_1,\eta_2\in\mathcal P_\mu(X),
\]
then \(\eta_1=\eta_2=\eta\). The intended ergodicity convention requires confirmation.

With these qualifications, the supplied questions are:
\begin{enumerate}
\item Describe the orbit closures for initial points \(x_0\in X\), resolving the source's notation \(\Gamma x_0\).
\item Determine whether the intended empirical measures \(\nu_n\) converge, specify the topology and probabilistic mode of convergence, and identify their limits. The applicable quantifier on \(x_0\) is not supplied.
\item Establish an appropriate no-escape-of-mass statement. For example, in a locally compact setting, one possible formulation for a deterministic sequence is tightness: for every \(\varepsilon>0\), there exists a compact \(K\subseteq X\) such that \(\nu_n(K)\geq1-\varepsilon\) for every \(n\geq1\). For random empirical measures, the probabilistic quantifiers and dependence of \(K\) must also be specified. This example is not an additional hypothesis or an asserted conclusion of the source.
\item Describe all \(\mu\)-stationary measures on \(X\), clarifying whether only probability measures are intended.
\item Describe the topology of the set of \(\mu\)-ergodic \(\mu\)-stationary measures, after specifying both ergodicity and the ambient topology on measures.
\end{enumerate}
No unconditional convergence or no-escape assertion follows from the supplied formulation. The background refers to results in particular homogeneous settings, but does not specify those settings or their hypotheses for this programme.
\subsection{Short English statement}
The record asks for orbit-closure classifications, empirical-measure convergence and limits, absence of mass loss, classification of stationary measures, and the topology of ergodic stationary measures. The acting objects, hypotheses, empirical measures, and convergence conventions remain unspecified, so this is a qualified research programme requiring human mathematical review.
\subsection{Sources}
\begin{itemize}
\item Supplied local UnsolvedMath registry record, permanent ID \(30003985\), exact source problem code \texttt{catalogue entry}, title \emph{Orbit Closures and Stationary Measures of Random Walks}, in the collection \emph{Oberwolfach Reports Open Problems}. No individual proposer is supplied; the proposed year is recorded as 2018.
\item Supplied report citation: \emph{Arbeitsgemeinschaft: Rigidity of Stationary Measure} (2019), Oberwolfach Reports, pp.~2805--2858, DOI \texttt{10.4171/owr/2018/47}. Source URL and supplied literature-source link, labelled \emph{Oberwolfach Report 47/2018}: \url{https://doi.org/10.4171/owr/2018/47}. The relationship between the citation year 2019 and the report designation 47/2018 is not explained in the record. No problem-specific page or section is supplied.
\item Literature assessment dated exactly \texttt{2026-08-21}: the supplied curation labels the problem open and states that the report proves orbit-closure, stationary-measure, and empirical-convergence results in particular homogeneous settings, while the broad five questions have no general resolution. This is supplied dated curation, not an independently verified literature assessment. The background explicitly requests independent verification of time-sensitive status and cited claims.
\item The registry labels the statement \texttt{corrected\_verified} and reports a manual semantic, typographic, and LaTeX review asserting self-containment and preservation of the target. These are supplied editorial claims; the formulation gaps identified here remain unresolved.
\end{itemize}
\end{document}
